\parte{Para Consultar}{Glossário, tabela de blocos e dicas para consertar bugs.}

\chapter*{Glossário}
\addcontentsline{toc}{chapter}{Glossário}

\begin{description}[leftmargin=!,labelwidth=3.6cm,style=nextline,itemsep=6pt]
\item[\textbf{Ator}] Personagem ou objeto do projeto (em inglês, \textit{sprite}).
\item[\textbf{Palco}] Área onde a ação acontece, com o cenário de fundo.
\item[\textbf{Bloco}] Peça de código que se encaixa em outras.
\item[\textbf{Bug}] Erro no programa. Consertá-lo se chama \textbf{depurar}.
\item[\textbf{Cenário}] Imagem de fundo do palco.
\item[\textbf{Clone}] Cópia temporária de um ator.
\item[\textbf{Condição}] Pergunta com resposta sim ou não, que o programa usa para decidir.
\item[\textbf{Coordenadas}] Os números \textbf{x} e \textbf{y} que dizem onde algo está no palco.
\item[\textbf{Evento}] Algo que acontece e dispara o código (clique, tecla, mensagem\ldots).
\item[\textbf{Fantasia}] Cada aparência de um ator.
\item[\textbf{Função}] Bloco criado por você (Meus Blocos).
\item[\textbf{Laço}] Repetição de um conjunto de blocos (\textbf{repita}, \textbf{sempre}).
\item[\textbf{Lista}] Conjunto de valores guardados em ordem.
\item[\textbf{Mensagem}] Aviso enviado de um ator para os outros.
\item[\textbf{Programa}] Conjunto de ordens dadas ao computador.
\item[\textbf{Roteiro}] Plano escrito do projeto antes de programar.
\item[\textbf{Variável}] Caixinha com nome que guarda um valor que pode mudar.
\end{description}

\chapter*{Tabela de Blocos}
\addcontentsline{toc}{chapter}{Tabela de Blocos}

\renewcommand{\arraystretch}{1.7}
\begin{center}
\begin{tabular}{@{}ll@{}}
\bloco[scMotion]{Movimento} & mova, gire, vá para, deslize, mude x/y, aponte \\
\bloco[scLooks]{Aparência} & diga, pense, fantasia, tamanho, efeitos, mostre/esconda \\
\bloco[scSound]{Som} & toque o som, volume, notas e tambores \\
\bloco[scEvents]{Eventos} & bandeira, teclas, cliques, mensagens \\
\bloco[scControl]{Controle} & espere, repita, sempre, se, senão, clones, pare \\
\bloco[scSensing]{Sensores} & tocando em, tecla, mouse, pergunte, cronômetro \\
\bloco[scOperators]{Operadores} & contas, comparações, e/ou/não, aleatório, junte \\
\bloco[scVariables]{Variáveis} & criar, defina, mude, mostre, esconda \\
\bloco[scLists]{Listas} & adicione, item, tamanho, apague \\
\bloco[scMyBlocks]{Meus Blocos} & blocos criados por você \\
\end{tabular}
\end{center}

\chapter*{Bugs Comuns}
\addcontentsline{toc}{chapter}{Bugs Comuns}

\begin{cuidado}[O ator não se mexe]
Os blocos estão soltos? Falta o bloco de evento no topo? Você está com o ator certo selecionado?
\end{cuidado}

\begin{cuidado}[O ator anda de cabeça para baixo]
Mude o \bloco[scMotion]{estilo de rotação} para esquerda-direita.
\end{cuidado}

\begin{cuidado}[A pontuação continua do jogo anterior]
Falta um \bloco[scVariables]{defina \menu{pontos} para \entrada{0}} no começo do projeto.
\end{cuidado}

\begin{cuidado}[A mensagem não chega]
O nome da mensagem está escrito diferente no \bloco[scEvents]{envie} e no \bloco[scEvents]{quando eu receber}? Escolha sempre no menu.
\end{cuidado}

\begin{cuidado}[O projeto ficou lento]
Você criou clones sem apagá-los? Use \bloco[scControl]{apague este clone} quando eles não forem mais necessários.
\end{cuidado}

\begin{cuidado}[A cor não é reconhecida]
Use o conta-gotas do bloco \bloco[scSensing]{tocando na cor} para escolher exatamente a mesma cor do desenho.
\end{cuidado}
